%document
\documentclass[10pt]{beamer}
%theme
\usetheme{metropolis}
% packages
\usepackage{color}
\usepackage{listings}
\usepackage[main=ngerman, USenglish]{babel}
\usepackage[utf8]{inputenc}
\usepackage{multicol}
\usepackage{csquotes}


% color definitions
\definecolor{mygreen}{rgb}{0,0.6,0}
\definecolor{mygray}{rgb}{0.5,0.5,0.5}
\definecolor{mymauve}{rgb}{0.58,0,0.82}
\definecolor{paleorange}{HTML}{CC7832}

\lstset{
    backgroundcolor=\color{white},
    % choose the background color;
    % you must add \usepackage{color} or \usepackage{xcolor}
    basicstyle=\footnotesize\ttfamily,
    % the size of the fonts that are used for the code
    breakatwhitespace=false,
    % sets if automatic breaks should only happen at whitespace
    breaklines=true,                 % sets automatic line breaking
    captionpos=b,                    % sets the caption-position to bottom
    commentstyle=\color{mygreen},    % comment style
    % deletekeywords={...},
    % if you want to delete keywords from the given language
    extendedchars=true,
    % lets you use non-ASCII characters;
    % for 8-bits encodings only, does not work with UTF-8
    frame=single,                    % adds a frame around the code
    keepspaces=true,
    % keeps spaces in text,
    % useful for keeping indentation of code
    % (possibly needs columns=flexible)
    keywordstyle=\color{blue},       % keyword style
    % morekeywords={*,...},
    % if you want to add more keywords to the set
    numbers=left,
    % where to put the line-numbers; possible values are (none, left, right)
    numbersep=5pt,
    % how far the line-numbers are from the code
    numberstyle=\tiny\color{mygray},
    % the style that is used for the line-numbers
    rulecolor=\color{black},
    % if not set, the frame-color may be changed on line-breaks
    % within not-black text (e.g. comments (green here))
    stepnumber=1,
    % the step between two line-numbers.
    % If it's 1, each line will be numbered
    stringstyle=\color{mymauve},     % string literal style
    tabsize=4,                       % sets default tabsize to 4 spaces
    % show the filename of files included with \lstinputlisting;
    % also try caption instead of title
    language = Java,
	showspaces = false,
	showtabs = false,
	showstringspaces = false,
	escapechar = ,
        morecomment=[s][\textcolor{paleorange}]{@}{\ },
}

\def\ContinueLineNumber{\lstset{firstnumber=last}}
\def\StartLineAt#1{\lstset{firstnumber=#1}}
\let\numberLineAt\StartLineAt



\newcommand{\codeline}[1]{
        \alert{\texttt{#1}}
}

\setmonofont{Fira Code}

% This Document contains the information about this course.

% Authors of the slides
\author{Marcus Köhler}

% Name of the Course
\institute{Java-Kurs}

% Fancy Logo
\titlegraphic{\hfill\includegraphics[height=1.25cm]{../templates/fsr_logo_cropped}}

\usepackage{ulem}
\usepackage{mathtools}

\title{Java}
\subtitle{Prozedurale Programmierung \& OOP I}
\date{\today}

\begin{document}

\begin{frame}
	\titlepage
\end{frame}
\begin{frame}{Überblick}
	\setbeamertemplate{section in toc}[sections numbered]
	\tableofcontents
\end{frame}

\section{Nachtrag: Arrays}
\begin{frame}[fragile]{Arrays}
    Ein \textit{Array}(oder auch \enquote{Feld}) ist eine Menge fester Größe von gleichartigen Datentypen bzw. Objekten.
    \begin{lstlisting}[gobble=8]
        int[] i_arr = {1, 2, 4, 8};
        char[] c_arr = {'c', 'h', 'a', 'r'};
        float[] f_arr = new float[5]; //empty array
        String[] s_arr; //Declaration
        s_arr = new String[10]; //Initialization
    \end{lstlisting}
\end{frame}

\begin{frame}[fragile]{Arrays}
    Die Elemente eines \textit{Arrays} sind mit einem Index versehen, über den man auf sie zugreifen kann; \\
    Java-Arrays sind 0-indexed(sprich: \enquote{zero-indexed}), d.h. das erste Element hat den Index 0.
    \begin{lstlisting}[gobble=8]
        System.out.println(i_arr[0]); //prints '1'
        System.out.println(c_arr[3]); //prints 'r'
    \end{lstlisting}
\end{frame}

\section{Prozedurale Programmierung}
\begin{frame}{Kontrollstatements}
	
    \begin{itemize}[<+->]
        \item Verzweigung \\
            \onslide<+->{\texttt{if, else, else if}} \\
            \onslide<+->{\texttt{switch-case}}
        \item Schleifen \\
            \onslide<+->{\texttt{while, do-while}} \\
            \onslide<+->{\texttt{for, for-each}}
	\end{itemize}
\end{frame}

\subsection{If-Else}
\begin{frame}[fragile]{if \& else}
    \begin{lstlisting}[gobble=8]
        if(condition) {
            // do something if condition is true
        } else if(another condition){
            // do if "else if" condition is true 
        } else {
            // otherwise do this
        }
    \end{lstlisting}
\end{frame}

\begin{frame}[fragile]{if \& else}
    \begin{lstlisting}[gobble=8]
        public class IfElseExample {
            public static void main(String[] args) {
                int myNumber = 5;
                
                if(myNumber == 3) {
                    System.out.println("Strange number");
                } else if(myNumber == 2) {
                    System.out.println("Unreachable code");
                } else {
                    System.out.println("Will be printed");
                }
            }
        }
    \end{lstlisting}
\end{frame}

\begin{frame}{Bedingungen?}
    Bedingungen sind Statements, welche einen \texttt{boolean}(Wahrheitswert) darstellen.
    Logische Verknüpfungen werden folgendermaßen dargestellt:
    \begin{itemize}[<+->]
        \item \texttt{!} Negation
        \item \texttt{==} Gleichheit
        \item \texttt{!=} Ungleichheit
        \item \texttt{>} Größer
        \item \texttt{>=} Größer/Gleich
        \item \texttt{<} Kleiner
        \item \texttt{<=} Kleiner/Gleich
    \end{itemize}
    \onslide<+->{\textit{Anmerkung}: Mehrere Bedingungen können mit \texttt{\&\&} (logisches UND)\\ oder \texttt{||} (logisches ODER) verkettet werden.
    Die Reihenfolge ist hierbei in einigen Fällen wichtig!}
\end{frame}

\subsection{while}
\begin{frame}[fragile]{while}
    \begin{lstlisting}[gobble=8]
        while(condition) {
            // do code while condition is true
        }
    \end{lstlisting}
\end{frame}

\begin{frame}[fragile]{while Beispiel}
    \begin{lstlisting}[gobble=8]
        public class WhileExample {

            public static void main(String[] args) {
                int a = 0;
                while(a <= 10) {
                    System.out.println(a);
                    a++; // Otherwise you would get an endless loop
                }
            }        
        }
    \end{lstlisting}
\end{frame}

\subsection{for}
\begin{frame}[fragile]{for}
    \begin{lstlisting}[gobble=8]
        for(initial value, condition, change) {
	        // do code while condition is true
        }
    \end{lstlisting}
\end{frame}

\begin{frame}[fragile]{for Beispiel}
    \begin{lstlisting}[gobble=8]
        public class ForExample {

            public static void main(String[] args) {
                for(int i = 0; i <= 10; i++) {
                    System.out.print("na ");
                }
                System.out.println("BATMAN!");
            }        
        }
    \end{lstlisting}
\end{frame}

\begin{frame}[fragile]{for}
    Das vorangegangene Programm kann man auch mit einer \texttt{while}-Schleife umsetzen:
    \begin{lstlisting}[gobble=8]
        public class ForWhileEx{
        public static void main(String[] args) {
            int i = 0;
            while(i <= 10) {
                System.out.println("na ");
                i++;
                }
            System.out.println("BATMAN!");
            }
        }
    \end{lstlisting}
\end{frame}

\subsection{Funktionen}
\begin{frame}[fragile]{Funktionen}
    \textit{Funktionen} sind sozusagen mehrmals nutzbare Code-Snippets:
    \begin{semiverbatim}
        \alert<2>{public int sq(int n)} \alert<3>{\{
            return n*n;
        \}}
    \end{semiverbatim}
    \onslide<2->{Funktionen bestehen aus einer \textit{Signatur}}
    \onslide<3->{und einem \textit{Körper}. \\}
    \onslide<4->{Im OOP werden Funktionen als \textit{Methoden} bezeichnet.}
\end{frame}

\begin{frame}[fragile]{Funktionssignaturen}
    \begin{center} 
    $
    {\color{orange}\overbracket[1pt]{\texttt{\large{public static}}}^{\text{access\ modifiers}}\ }
    {\color{blue}\underbracket[1pt]{\texttt{\large{void}}}_{\text{return\ type}}\ }
    {\color{black}\overbracket[1pt]{\texttt{\large{main}}}^{\text{name}}\ }
    {\color{gray}\underbracket[1pt]{\texttt{\large{(String[] args)}}}_{\text{parameter(s)}}}
    $
    \end{center}
    \onslide<+-> Die \texttt{access modifier} regeln, wer die Funktion aufrufen darf. \\
    \onslide<+-> Der \texttt{return type}(auch: Rückgabewert) gibt an, welchen Typ das \enquote{Ergebnis} hat.\\
    \onslide<+-> Die \texttt{Parameter}(auch: Argumente) sind die Daten, die die Methode braucht bzw. erwartet.
\end{frame}

\section{OOP in Java}

\begin{frame}{}
	\begin{center}
		{\huge Objektorientierte Programmierung}
	\end{center}
\end{frame}

\subsection{Grundlagen}

\begin{frame}[fragile]{Class Student}
    \begin{lstlisting}[gobble=8]
        public class Student {

            // Attributes
            private String name; 
            private int matriculationNumber; 

            // Methods
            public void setName(String name) {
                this.name = name;
            }

            public int getMatriculationNumber() {
                return matriculationNumber;
            }
        }
    \end{lstlisting}

% What is visible here:
% Attributes store the state of the object
% Methods implement the behaviour of the object

\end{frame}

\begin{frame}[fragile]{Ein Objekt erstellen}
	Ein primitiver Datentyp wird so erstellt:
    \begin{lstlisting}[gobble=8]
	    int a; // declare a
	    a = 273; // assign 273 to a
	\end{lstlisting}
    \pause
	Bei Objekten funktioniert das ähnlich:
    \begin{lstlisting}[gobble=8]
	    Student example = new Student(); 
		// create an instance of Student
	\end{lstlisting}
    \pause
    Ein Objekt einer \textit{Klasse} wird auch \textit{Instanz} genannt.
	Eine Variable, die mit einem Objekt belegt ist, heißt \textit{Referenz}.
\end{frame}

\subsection{Methoden}
\begin{frame}[fragile]{Eine Methode aufrufen}
	\begin{lstlisting}[gobble=8]
        public class Student {
        
            private String name;
        
            public String getName() {
                return name;
            }
            
            public void setName(String newName) {
                name = newName;
            }
        }
	\end{lstlisting}
    Die Klasse \emph{Student} (und damit das Objekt) hat zwei Methoden: \\  \texttt{void printTimetable()} und \texttt{void printName()}.
\end{frame}

\begin{frame}[fragile]{Eine Methode aufrufen}
    \begin{lstlisting}[gobble=8]
        public class Main {

            public static void main(String[] args) {
                Student example = new Student(); // creation
                example.setName("Jane"); // method call
                String name = example.getName(); 
                System.out.println(name); // Prints "Jane"
            } 
        }
    \end{lstlisting}
	Man kann eine Methode eines Objekts nach der Instanzierung mittels \texttt{Referenz.methodName();} aufrufen.
\end{frame}

\begin{frame}[fragile]{Eine Methode aufrufen}
    \begin{lstlisting}[gobble=8]
        public class Student {
        
            private String name;
        
            public void setName(String newName) {
                name = newName;
                printName();   // Call own method
                this.printName(); // Or this way
            }
            
            public void printName() {
                System.out.println(name);
            }
        }
    \end{lstlisting}
	Ein Objekt kann seine eigenen Methoden mit \texttt{methodName();} oder \texttt{this.methodName();} aufrufen.
\end{frame}

\begin{frame}[fragile]{Methoden mit Parametern}
    \begin{lstlisting}[gobble=8]
        public class Calc {

            public void add(int summand1, int summand2) {
                System.out.println(summand1 + summand2);
            }
                
            public static void main(String[] args) {
                int summandA = 1;
                int summandB = 2;
                Calc calculator = new Calc();
                System.out.print("1 + 2 = ");
                calculator.add(summandA, summandB); 
                // prints: 3
            }    
        }
	\end{lstlisting}
\end{frame}

\subsection{Rückgabewerte}
\begin{frame}[fragile]{Methoden mit Rückgabewerten}
	Eine Methode ohne Rückgabewert hat den Typ \texttt{void}:
	\begin{lstlisting}[gobble=8]
        public void add(int summand1, int summand2) {
            System.out.println(summand1 + summand2);
        }
	\end{lstlisting}
    Methoden mit (passenden) Rückgabewerten können Variablen zugewiesen werden:
    \begin{lstlisting}[gobble=8]
        int a, b, aTimesB;
        int a = 5;
        int b = 7; 
        int aTimesB = mul(a, b);
    \end{lstlisting}
\end{frame}

\subsection{Der Konstruktor}
\begin{frame}[fragile]{Konstruktoren}
    Ein \textit{Konstruktor} ist eine spezielle Methode, welche den \enquote{Urzustand} eines Objekts definiert.\\
    Er wird aufgerufen, sobald ein Objekt instanziert wird:
    \begin{lstlisting}[gobble=8]
        public class Calc {
            private int summand1;
            private int summand2;
        
            public Calc() { // Constructor
                summand1 = 0;
                summand2 = 0;
            }    

            public static void main(String[] args) {
                Calc calc = new Calc(); //Constructor is called
            }
        }
	\end{lstlisting}
\end{frame}

\begin{frame}[fragile]{Konstruktoren mit Parametern}
    Ein Konstruktor kann, wie eine \enquote{normale} Methode, auch Parameter übergeben bekommen. \\
    Das ist nötig, wenn Anfangswerte etc. angegeben werden:
    \begin{lstlisting}[gobble=8]
        public class Calc {
            private int summand1;
            private int summand2;
        
            public Calc(int x, int y) {
                summand1 = x;
                summand2 = y;
            }
        }
	\end{lstlisting}
    \begin{lstlisting}[gobble=8]
	    [...]
	    Calc myCalc = new Calc(7, 9);
	\end{lstlisting} 
\end{frame}

\begin{frame}{Ausblick}
    \begin{itemize}
        \item Datenkapselung
        \item Vererbung
        \item Polymorphie
        \item Das \texttt{this} Keyword
    \end{itemize}
\end{frame}
\end{document}
